\section{Imitation and Self-Preferencing}

\begin{frame}
  \frametitle{What Is Added?}
  \footnotesize
  This part adds two practices to the baseline dual-mode model: \textbf{product imitation} and \textbf{Self-preferencing}.

  \begin{block}{Product imitation}
    \begin{itemize}
      \item \textbf{Assumption:} copied product gives surplus $\nu+\Delta$.
      \item \textbf{Choice:} $S$: participate or not; then $M$: imitate or not
        if $S$ participates.
      \item \textbf{Motivation:} marketplace data helps $M$ identify and copy
        successful innovations.
    \end{itemize}
  \end{block}

  \begin{block}{Self-preferencing}
    \begin{itemize}
      \item \textbf{Assumption:} consumers discover $S$ through $M$'s recommendation.
      \item \textbf{Choice:} $M$: recommend $S$ or not after prices are set.
      \item \textbf{Motivation:} consumers often rely on platform search, ranking,
        or recommendations to find new products.
    \end{itemize}
  \end{block}

  % \begin{alertblock}{Purpose}
  %   Test a dual-mode ban under perfect imitation and perfect steering.
  % \end{alertblock}
\end{frame}

\begin{frame}
  \frametitle{Timing: Where Section 4 Changes the Game}
  \footnotesize
  \begin{center}
    \begin{tikzpicture}[
      >=stealth,
      stage/.style={
        draw=black!35,
        fill=black!3,
        minimum height=1.35cm,
        text width=2.75cm,
        align=center,
        inner sep=4pt
      },
      changed/.style={
        draw=darkred!85!black,
        fill=darkred!7,
        very thick,
        minimum height=1.35cm,
        text width=2.75cm,
        align=center,
        inner sep=4pt
      },
      note/.style={
        draw=darkred!85!black,
        fill=darkred!5,
        text width=2.85cm,
        align=center,
        inner sep=3pt
      }
    ]
      \node[stage] (s1) at (0,0)
        {\textbf{Stage 1}\\ $M$: mode, $\tau$};
      \node[changed] (s2) at (3.55,0)
        {\textbf{Stage 2}\\ $S$: join, $\Delta$\\ $M$: imitate?};
      \node[stage] (s3) at (7.10,0)
        {\textbf{Stage 3}\\ prices\\ $p_m,p_i,p_o$};
      \node[changed] (s4) at (10.65,0)
        {\textbf{Stage 4}\\ $M$: show $S$?\\ consumers buy};

      \draw[->,thick] (s1) -- (s2);
      \draw[->,thick] (s2) -- (s3);
      \draw[->,thick] (s3) -- (s4);

      \node[note] (n2) at (3.55,-1.95)
        {\textbf{New: imitation}\\ only if $S$ participates and $M$ is in dual mode};
      \node[note] (n4) at (10.65,-1.95)
        {\textbf{New: steering}\\ discovery depends on $M$'s recommendation};

      \draw[->,darkred!85!black,thick] (n2.north) -- (s2.south);
      \draw[->,darkred!85!black,thick] (n4.north) -- (s4.south);
    \end{tikzpicture}
  \end{center}

  \textbf{Changes from the baseline}
  \begin{itemize}
    \item $M$ copying choice: $M$ imitates or not.
    \item $M$ disclosure choice: $M$ shows $S$ or not.
  \end{itemize}
\end{frame}

\begin{frame}
  \frametitle{Marketplace and Seller Mode}
  \footnotesize
  Marketplace mode: same commission $\tau$ from any marketplace sale.

  \begin{columns}[T,onlytextwidth]
    \begin{column}{0.48\textwidth}
      \begin{block}{If $M$ shows $S$}
        Buy $S$ through the marketplace iff:
        \begin{equation}
          \Delta+b-p_i\geq \max\{\Delta-p_o,\ b-\tau,\ 0\}
          \label{eq:show-s-marketplace}
        \end{equation}
        \begin{itemize}
          \item Left: surplus from $S$ on $M$.
          \item Right: other best outside option.
        \end{itemize}
      \end{block}
    \end{column}

    \begin{column}{0.48\textwidth}
      \begin{block}{If $M$ hides $S$}
        Buy a fringe seller on the marketplace iff:
        \begin{equation}
          b-\tau\geq 0
          \label{eq:hide-s-fringe}
        \end{equation}
        \begin{itemize}
          \item Left: surplus from a fringe seller on $M$.
          \item Right: no purchase.
        \end{itemize}
      \end{block}
    \end{column}
  \end{columns}

  \begin{block}{Recommendation rule}
    \begin{itemize}
      \item Show $S$ if Eq. \eqref{eq:show-s-marketplace}, holds: $M$ earns commission $\tau$.
      \item Hide $S$ if Eq. \eqref{eq:show-s-marketplace}, fails but Eq. \eqref{eq:hide-s-fringe}, holds.
    \end{itemize}
  \end{block}

  Seller mode is unchanged: $M$ sells only its own product.
\end{frame}


% \begin{frame}
%   \frametitle{Seller mode}

%   Seller mode: $M$ sells only its own product; no marketplace listing for $S$ (unchanged vs baseline).

%   \begin{block}{Baseline expression}
%     Constraint only from fringe sellers' direct channel:
%     \begin{equation}
%       \Pi^{sell} = \max_{p_m\leq b+\sigma} \; p_m G(\nu+b+\sigma-p_m).
%     \end{equation}
%   \end{block}

%   {Same equilibrium implication}
%     \begin{itemize}
%       \item Constraint binds: $p_m^*=b+\sigma$.
%       \item $M$ sells to all consumers: $\Pi^{sell}=(b+\sigma)G(\nu)$.
%       \item $S$ sells to no one: $\pi^{sell}=0$, $\Delta^{sell}=\Delta^l$.
%     \end{itemize}

% \end{frame}

\begin{frame}
  \frametitle{Dual Mode: Recommendation Decision}
  % \footnotesize

  \begin{block}{What does $M$ compare?}
    \begin{itemize}
      \item Own expected margin from selling $M$'s product or imitation.
      \item Expected commission from showing $S$.
    \end{itemize}
  \end{block}

  \begin{block}{Recommendation rule}
    \begin{itemize}
    \item $M$ choose to hide $S$
    \begin{itemize}
      \item $M$'s expected margin $>$ expected commission from $S$;
      \item Or the marketplace-purchase condition for $S$, Eq. \eqref{eq:show-s-marketplace}, fails.
    \end{itemize}
    \item $M$ choose to show $S$ 
    \begin{itemize}
      \item $M$'s expected margin $<$ expected commission from $S$;
      \item and the marketplace-purchase condition for $S$, Eq. \eqref{eq:show-s-marketplace}, holds.
    \end{itemize}
  \end{itemize}
  \end{block}

\end{frame}



\begin{frame}
  \frametitle{Pricing Without Imitation}
  \footnotesize
  Without imitation, the subgame can be exploitative or price-squeezing.
  \begin{enumerate}
    \item \textbf{Exploitative}: $M$ hides $S$, and charges the highest price allowed by outside options.
    \item \textbf{Price squeeze}: $M$ shows $S$, but caps $S$'s inside price and margin.
  \end{enumerate}

  \begin{block}{Exploitative equilibrium}

    \begin{itemize}
      \item \textbf{Exploitative price}: 
      \begin{equation}
        p_m^*=\min\{\tau,b\}+\sigma
        \label{eq:pm-wo-imitation}
      \end{equation}
      \begin{itemize}
        \item $M$ vs fringe through marketplace: $\nu + \sigma + b - p_m^* > \nu + b - \tau$ $\Longrightarrow$ $p_m^* < \tau + \sigma$
        \item $M$ vs fringe not through marketplace: $\nu + \sigma + b - p_m^* > \nu $ $\Longrightarrow$ $p_m^* < b + \sigma$
      \end{itemize}
      \item \textbf{Exploitative profit}: using the seller-mode profit function, Eq. \eqref{eq:seller-profit}, and Eq. \eqref{eq:pm-wo-imitation}:
      \begin{equation}
        \begin{aligned}
          \Pi_{no-imi}^{exploit} &= p^*_m \times G(\nu+b+\sigma-p^*_m) \\
          & = (\min\{\tau,b\}+\sigma)G(\nu+b-\min\{\tau,b\}).
        \end{aligned}
      \end{equation}
    \end{itemize}
  \end{block}

  
\end{frame}

\begin{frame}
  \frametitle{Pricing Without Imitation: Price Squeeze}
  \footnotesize

  Without imitation, products and quality gaps are the same as in the baseline.
  The only new instrument is steering:
  \begin{itemize}
    \item If $S$ is shown, pricing follows the baseline price-squeeze logic.
    \item Self-preferencing adds one more constraint: $M$ must prefer showing $S$ to hiding it.
  \end{itemize}


\end{frame}

\begin{frame}
  \frametitle{Pricing With Imitation \footnote{Imitation matters when $S$'s product is better than $M$'s original product $\Delta>\sigma$.}}
  \footnotesize
  $M$ can sells the \textbf{copied product ($\sigma = \Delta$)} to all consumers at
  \begin{block}{Exploitative equilibrium}
    \begin{equation}
      \begin{aligned}
        p_m^*&=\min\{\tau,b\}+ \sigma \qquad (\text{imitation: }\sigma = \Delta) \\
        & = \min\{\tau,b\}+ \textcolor{red}{\Delta} .
      \end{aligned}
    \end{equation}

    \begin{equation}
      \begin{aligned}
        \tilde{\Pi}^{exploit}_{imi} & = p^*_m \times G(\nu+b+\textcolor{red}{\Delta}-p^*_m) \\
        & = (\min\{\tau,b\}+\Delta)G(\nu+b-\min\{\tau,b\}).
      \end{aligned}
    \end{equation}
  \end{block}

  \begin{block}{Price-squeeze equilibrium}
    \begin{equation}
      \begin{aligned}
        p_i^* &=p_m^* + \Delta - \sigma \qquad (\text{imitation: }\sigma = \Delta) \\
        & = \min\{\tau,b\} = \tau .
      \end{aligned}
    \end{equation}

    \textbf{notice:} if $M$ can imitation $p_i^* = \tau$ does not depend on $\Delta$, $S$ will not innovate.
  \end{block}

  % \begin{alertblock}{Mechanism}
  %   Imitation extracts innovation surplus; steering prevents $S$ from using
  %   discovery to discipline the platform.
  % \end{alertblock}
\end{frame}

\begin{frame}
  \frametitle{Proposition 5: Dual Mode Equilibrium}
  \scriptsize
  \begin{block}{Assumptions in the regulated-practice environment}
    \begin{itemize}
      \item \textbf{Imitation allow:} $M$ can perfectly imitate $S$'s innovation at no cost.
      \item \textbf{Self-preferencing allow:} $M$ can perfectly steer consumers toward its own offer.
      \item \textbf{Dual mode is feasible:} $M$ both sells its own product and hosts $S$.
    \end{itemize}
  \end{block}

  \begin{block}{Main variables in equilibrium}
    \begin{itemize}
      \item $M$'s profit: $\Pi^{dual}= p^*_m G(\nu+b+\Delta^{dual}-p^*_m) = (b+\max\{\sigma,\Delta^l\})G(\nu).$
      \item $S$'s operating profit: $\pi^{dual}=(p_i^*-\tau)G(\nu+b+\Delta^{dual}-p_i^*)=0$.
      \item Innovation level: $\Delta^{dual} = \Delta^l$ $\quad \leftarrow \quad$ $\Delta^{dual}= \arg\max_{\Delta^{dual} \geq \Delta^l} \pi^{dual} - K(\Delta)$.
      \item Consumer surplus: $CS^{dual}=CS^0\equiv \int_0^\infty \max\{\nu,\nu_o\}\,dG(\nu_o)$.
      \item Total welfare: $W^{dual}=\Pi^{dual}+\pi^{dual}+CS^{dual}$
    \end{itemize}
  \end{block}
\end{frame}

\begin{frame}
  \frametitle{Proposition 5: Outcome Vector}
  \scriptsize
 
  \begin{center}
    \textbf{Table: Dual-mode benchmark with imitation and self-preferencing}

    \vspace{0.2em}
    \setlength{\tabcolsep}{1.6pt}
    \begin{tabular}{@{}p{0.14\linewidth}p{0.24\linewidth}p{0.17\linewidth}cccc@{}}
      \toprule
      Case & $M$'s choice of mode & $\Pi$ & $\pi$ & $\Delta$ & $CS$ & $W$ \\
      \midrule
      $\sigma\geq\Delta^l$
        & Dual - no imitation, $M$ sells
        & $(b+\sigma)G(\nu)$
        & $0$
        & $\Delta^l$
        & $CS^0$
        & $\Pi+CS^0$ \\
      $0<\sigma<\Delta^l$
        & Dual - high fee or imitation
        & $(b+\Delta^l)G(\nu)$
        & $0$
        & $\Delta^l$
        & $CS^0$
        & $\Pi+CS^0$ \\
      $\sigma\leq0$
        & Dual - high fee or imitation
        & $(b+\Delta^l)G(\nu)$
        & $0$
        & $\Delta^l$
        & $CS^0$
        & $\Pi+CS^0$ \\
      \bottomrule
    \end{tabular}
  \end{center}

  \begin{block}{High-fee and imitation paths}
    \begin{itemize}
      \item \textbf{High fee} ($M$ as marketplace): $\tau = b + \Delta^l$.
      \item \textbf{Imitation} ($M$ as seller): $p_m = b + \Delta^l$.
    \end{itemize}
  \end{block}

  \begin{itemize}
    \item $\pi=0$: with a high fee, $M$ extracts all of $S$'s price margin;
      with imitation, $M$ makes all sales itself.
    \item $\Delta=\Delta^l$: since $S$'s operating profit does not increase
      with $\Delta$, it has no incentive to innovate above the minimum level.
  \end{itemize}

\end{frame}

\begin{frame}
  \frametitle{Policy Map: How the Bans Are Related}
  \scriptsize

  \begin{block}{Ban types}
    \begin{itemize}
    \item \textbf{Structural ban}: change $M$'s business model by removing dual
      mode; it can shift $M$ to a pure seller or pure marketplace mode.
    \item \textbf{Behavioral ban}: keep dual mode possible but remove specific
      practices; it can restore price discipline and/or innovation incentives.
  \end{itemize}
  \end{block}

  \begin{center}
    \begin{tikzpicture}[
      >=stealth,
      base/.style={
        draw=black!45,
        fill=black!3,
        rounded corners=2pt,
        text width=3.05cm,
        align=center,
        inner sep=3pt
      },
      policy/.style={
        draw=darkred!85!black,
        fill=darkred!7,
        rounded corners=2pt,
        text width=3.05cm,
        align=center,
        inner sep=3pt
      },
      result/.style={
        draw=black!40,
        fill=black!2,
        rounded corners=2pt,
        text width=2.05cm,
        align=center,
        inner sep=2pt
      }
    ]
      \node[base] (bench) at (0,1.25)
        {\textbf{Proposition 5}\\ dual mode with imitation\\ and self-preferencing};
      \node[policy] (struct) at (-3.55,-0.25)
        {\textbf{Structural ban}\\ remove dual mode from\\ $M$'s choice set};
      \node[policy] (behav) at (3.55,-0.25)
        {\textbf{Behavioral ban}\\ keep dual mode possible,\\ restrict practices};

      \node[result] (p6) at (-3.55,-1.75)
        {\textbf{Prop. 6}\\ dual-mode ban};
      \node[result] (p7) at (0.55,-1.75)
        {\textbf{Prop. 7}\\ no self-pref.};
      \node[result] (p8) at (3.55,-1.75)
        {\textbf{Prop. 8}\\ no imitation};
      \node[result] (p9) at (6.55,-1.75)
        {\textbf{Prop. 9}\\ neither practice};

      \draw[->,thick] (bench.south west) -- (struct.north);
      \draw[->,thick] (bench.south east) -- (behav.north);
      \draw[->,thick] (struct) -- (p6);
      \draw[->,thick] (behav.south) -- (p7.north);
      \draw[->,thick] (behav.south) -- (p8.north);
      \draw[->,thick] (behav.south) -- (p9.north);
    \end{tikzpicture}
  \end{center}
  
\end{frame}

\begin{frame}
  \frametitle{Dual-Mode Ban -- Step 1: What Is Banned?}
  \small
  \begin{block}{Policy object}
    The ban removes the platform's ability to operate in dual mode within this
    product category.
  \end{block}

  \begin{itemize}
    \item \textbf{Economic meaning}: $M$ cannot both host $S$ on its marketplace
      and sell its own competing product in the same category.
    \item \textbf{What it does not directly ban}: product imitation or
      self-preferencing as practices.
    \item \textbf{Mathematical meaning}: remove dual mode from $M$'s stage-0
      choice set,
      \[
        \{\text{Marketplace},\text{Seller},\text{Dual}\}
        \quad\longrightarrow\quad
        \{\text{Marketplace},\text{Seller}\}.
      \]
  \end{itemize}

  \begin{block}{Post-ban mode choice}
    \begin{equation}
      M \text{ chooses } \arg\max\{\Pi^{mkt},\Pi^{sell}\}.
    \end{equation}
  \end{block}
\end{frame}

\begin{frame}
  \frametitle{Dual-Mode Ban -- Computing the Post-Ban Outcome}
  \scriptsize
  Once dual mode is removed, $M$ compares the two pure modes.

  \begin{center}
    \setlength{\tabcolsep}{3pt}
    \begin{tabular}{@{}lccccc@{}}
      \toprule
      Mode & $\Pi$ & $\pi$ & $\Delta$ & $CS$ & $W$ \\
      \midrule
      Seller
        & $(b+\sigma)G(\nu)$
        & $0$
        & $\Delta^l$
        & $CS^0$
        & $(b+\sigma)G(\nu)+CS^0$ \\
      Marketplace
        & $bG(\nu)$
        & $\Delta^{mkt}G(\nu)-K(\Delta^{mkt})$
        & $\Delta^{mkt}$
        & $CS^0$
        & $(b+\Delta^{mkt})G(\nu)-K(\Delta^{mkt})+CS^0$ \\
      \bottomrule
    \end{tabular}
  \end{center}

  \begin{block}{Mode selected after the ban}
    If $\sigma>0$, $M$ chooses seller mode; if $\sigma\leq0$, $M$ chooses
    marketplace mode.
  \end{block}
\end{frame}

\begin{frame}
  \frametitle{Dual-Mode Ban -- Step 2: Equilibrium Result}
  \small
  \textbf{Proposition 6.} Compare the post-ban pure-mode outcome with the
  Proposition 5 dual-mode benchmark.

  \begin{center}
    \scriptsize
    \begin{tabular}{@{}lcccccc@{}}
      \toprule
      Region & $M$'s choice of mode & $\Pi$ & $\pi$ & $CS$ & $\Delta$ & $W$ \\
      \midrule
      $\sigma\geq\Delta^l$
        & Seller & $=$ & $=$ & $=$ & $=$ & $=$ \\
      $0<\sigma<\Delta^l$
        & Seller & $\downarrow$ & $=$ & $=$ & $=$ & $\downarrow$ \\
      $\sigma\leq0$
        & Marketplace & $\downarrow$ & $\uparrow$ & $=$ & $\uparrow$ & $\uparrow$ \\
      \bottomrule
    \end{tabular}
  \end{center}

  \begin{alertblock}{Interpretation}
    A structural ban changes $M$'s business model, but it does not directly
    restore on-platform competition or innovation incentives.
  \end{alertblock}
\end{frame}


\begin{frame}
  \frametitle{Alternative Policy: Behavioral Remedies}
  \small
  Instead of banning the dual role itself, the regulator can target specific
  conduct within dual mode.

  \begin{center}
    \scriptsize
    \begin{tabular}{@{}>{\raggedright\arraybackslash}p{0.28\linewidth}>{\raggedright\arraybackslash}p{0.34\linewidth}>{\raggedright\arraybackslash}p{0.30\linewidth}@{}}
      \toprule
      Remedy & What is banned & Main channel restored \\
      \midrule
      Ban self-preferencing only
        & $M$ must show $S$ when $S$ is listed
        & on-platform competition; showrooming constraint \\
      Ban imitation only
        & $M$ cannot copy $S$'s innovation
        & $S$'s innovation incentive, if commitment is valuable \\
      Ban both practices
        & no steering and no copying
        & baseline dual-mode mechanisms \\
      \bottomrule
    \end{tabular}
  \end{center}

  \begin{alertblock}{Logic}
    For each behavioral remedy, use the same two-step structure: identify the
    banned conduct, then compare the resulting equilibrium with Proposition 5.
  \end{alertblock}
\end{frame}

\begin{frame}
  \frametitle{Self-Preferencing Ban -- Step 1: What Is Banned?}
  \small
  \begin{block}{Policy object}
    $M$ must always disclose/show $S$'s product whenever $S$ participates on
    the marketplace.
  \end{block}

  \begin{itemize}
    \item \textbf{Economic meaning}: $M$ cannot steer consumers away from $S$;
      consumers observe all on-platform offers.
    \item \textbf{What remains allowed}: product imitation by $M$.
    \item \textbf{Mathematical meaning}: the showrooming constraint is restored,
      so $M$ cannot set a fee above the outside-channel discipline:
      \[
        \tau\leq b.
      \]
  \end{itemize}

  \begin{block}{Post-intervention dual-mode outcome}
    \begin{equation}
      \tau=b,\qquad \Delta=\Delta^l,\qquad
      \Pi=(b+\max\{\sigma-\Delta^l,0\})G(\nu+\Delta^l),\qquad \pi=0.
    \end{equation}
  \end{block}
\end{frame}

\begin{frame}
  \frametitle{Self-Preferencing Ban -- Step 2: Equilibrium Result}
  \small
  \textbf{Proposition 7.} Let $\underline{\sigma}^{steer}$ be the cutoff at
  which $M$ is indifferent between seller mode and dual mode without
  self-preferencing.

  \begin{center}
    \scriptsize
    \begin{tabular}{@{}lcccccc@{}}
      \toprule
      Region & $M$'s choice of mode & $\Pi$ & $\pi$ & $CS$ & $\Delta$ & $W$ \\
      \midrule
      $\sigma\geq\Delta^l$
        & Seller & $=$ & $=$ & $=$ & $=$ & $=$ \\
      $\underline{\sigma}^{steer}<\sigma<\Delta^l$
        & Seller & $\downarrow$ & $=$ & $=$ & $=$ & $\downarrow$ \\
      $\sigma\leq\underline{\sigma}^{steer}$
        & Dual, no self-pref. & $\downarrow$ & $=$ & $\uparrow$ & $=$ & $\uparrow$ \\
      \bottomrule
    \end{tabular}
  \end{center}

  \begin{alertblock}{Interpretation}
    This ban restores price competition, but not innovation: imitation still
    keeps $S$'s equilibrium profit at zero.
  \end{alertblock}
\end{frame}

\begin{frame}
  \frametitle{Imitation Ban -- Step 1: What Is Banned?}
  \small
  \begin{block}{Policy object}
    $M$ is prohibited from copying $S$'s product innovation.
  \end{block}

  \begin{itemize}
    \item \textbf{Economic meaning}: $M$ can credibly commit not to imitate,
      so \(S\)'s innovation may become valuable again.
    \item \textbf{What remains allowed}: self-preferencing/steering by $M$.
    \item \textbf{Mathematical meaning}: $M$ can choose a fee that respects the
      innovation constraint
      \[
        \tau\leq\bar{\tau},
      \]
      where $\bar{\tau}$ is the fee cutoff that keeps $S$ willing to choose
      high innovation.
  \end{itemize}
\end{frame}

\begin{frame}
  \frametitle{Imitation Ban -- Step 2: Equilibrium Result}
  \scriptsize
  \textbf{Proposition 8.} The ban has an effect only when it induces $M$ to set
  $\tau=\bar{\tau}$ and preserve $S$'s innovation incentive.

  \begin{block}{Condition for the innovation-preserving equilibrium}
    \begin{equation}
      b+\Delta^l\leq\bar{\tau}
      \quad\text{or}\quad
      \bar{\tau}G(\nu+\sigma+b)
      \geq (b+\max\{\sigma,\Delta^l\})G(\nu).
    \end{equation}
  \end{block}

  \begin{center}
    \scriptsize
    \begin{tabular}{@{}p{0.38\linewidth}cccccc@{}}
      \toprule
      Region & $M$'s mode & $\Pi$ & $\pi$ & $CS$ & $\Delta$ & $W$ \\
      \midrule
      $b+\Delta^l>\bar{\tau}$ and condition fails
        & Dual, no imitation & $=$ & $=$ & $=$ & $=$ & $=$ \\
      $b+\Delta^l\leq\bar{\tau}$ or condition holds
        & Dual, no imitation & $\uparrow$ & $=$ & $\uparrow$ & $\uparrow$ & $\uparrow$ \\
      \bottomrule
    \end{tabular}
  \end{center}
\end{frame}

\begin{frame}
  \frametitle{Ban Both Practices -- Step 1: What Is Banned?}
  \small
  \begin{block}{Policy object}
    $M$ is prohibited from both self-preferencing and product imitation.
  \end{block}

  \begin{itemize}
    \item \textbf{Economic meaning}: consumers see $S$'s product, and $S$ is
      protected from copying by $M$.
    \item \textbf{Mathematical meaning}: the model returns to the baseline
      dual-mode mechanism from Proposition 3.
    \item The showrooming constraint disciplines $\tau$, and the absence of
      imitation can restore $S$'s innovation incentive.
  \end{itemize}
\end{frame}

\begin{frame}
  \frametitle{Ban Both Practices -- Step 2: Equilibrium Result}
  \scriptsize
  \textbf{Proposition 9.} Let $\underline{\sigma}$ be the baseline cutoff at
  which $M$ switches between seller mode and dual mode after both practices are
  banned.

  \begin{center}
    \scriptsize
    \setlength{\tabcolsep}{2pt}
    \begin{tabular}{@{}p{0.20\linewidth}p{0.30\linewidth}ccccc@{}}
      \toprule
      Region & $M$'s choice of mode & $\Pi$ & $\pi$ & $CS$ & $\Delta$ & $W$ \\
      \midrule
      $\sigma\geq\max\{\underline{\sigma},\Delta^l\}$
        & Seller & $=$ & $=$ & $=$ & $=$ & $=$ \\
      $\underline{\sigma}<\sigma<\Delta^l$
        & Seller & $\downarrow$ & $=$ & $=$ & $=$ & $\downarrow$ \\
      $\sigma\leq\underline{\sigma}$
        & Dual, no imitation or self-pref.
        & $\downarrow$
        & $\uparrow^\dagger$ or $=$
        & $\uparrow$
        & $\uparrow^\ddagger$ or $=$
        & $\uparrow$ \\
      \bottomrule
    \end{tabular}
  \end{center}

  \begin{itemize}
    \item $\dagger$: $\pi$ increases if $b<\bar{\tau}$; otherwise it is unchanged.
    \item $\ddagger$: $\Delta$ increases if $b<\bar{\tau}$ or condition (8) holds; otherwise it is unchanged.
  \end{itemize}

  \begin{alertblock}{Interpretation}
    Banning both practices restores both channels: price discipline from
    visibility and innovation incentives from no imitation.
  \end{alertblock}
\end{frame}

\begin{frame}
  \frametitle{Policy Comparison}
  \small
  \begin{itemize}
    \item \textbf{Structural ban}: removes dual mode, but does not directly fix
      imitation or self-preferencing.
    \item \textbf{Ban self-preferencing}: restores price competition, but
      innovation remains low because imitation is still possible.
    \item \textbf{Ban imitation}: can restore innovation when $M$ benefits from
      committing not to copy.
    \item \textbf{Ban both}: closest to the baseline dual-mode mechanism and
      can restore both price competition and innovation.
  \end{itemize}

  \begin{alertblock}{Policy lesson}
    The harmful conduct is often a better regulatory target than the platform's
    dual role itself.
  \end{alertblock}
\end{frame}


\begin{frame}
  \frametitle{Section 4 Takeaways}
  \small
  \begin{enumerate}
    \item Perfect imitation and perfect steering make dual mode weakly best for $M$.
    \item In dual mode, $S$ chooses $\Delta^l$ because innovation rents are extracted.
    \item A dual-mode ban does not directly fix copying or steering.
    \item Banning self-preferencing restores price competition but not innovation.
    \item Banning imitation restores innovation when commitment is valuable enough.
    \item Banning both practices restores the baseline dual-mode mechanisms.
  \end{enumerate}

  \begin{alertblock}{Main policy message}
    The harmful conduct is often a better regulatory target than the platform's
    dual role itself.
  \end{alertblock}
\end{frame}
